
\setcounter{subsection}{3} \setcounter{subsubsection}{-1}

\subsubsection{}
\label{XIII.5.3.0} Retomemos las hipótesis y las notaciones
de~\Ref{XIII.5.2}, suponiendo además que
\ifthenelse{\boolean{orig}}
{$S$}
{$X$}
es estrictamente local. Entonces se deduce de \loccit que todo
recubrimiento étale conexo de~$U$, moderadamente ramificado
relativamente a~$D$, es cociente de un recubrimiento moderadamente
ramificado de la forma
$$
U'=U[T_1,\dots,T_r]/(T_1^{n_1}-f_1,\dots,T_r^{n_r}-f_r),
$$
donde los $n_i$ son enteros primos con~$p$. Sea $\bbmu_n$ el grupo de
las raíces $n$-ésimas de la unidad de~$U$. El grupo de los
$U$-automorfismos de~$U'$ no es otro que el grupo
$\bbmu_{n_1}\times\dots\times\bbmu_{n_r}$: una raíz $n_i$-ésima de
la unidad~$\xi_i$ actúa sobre~$U'$ transformando $T_i$ en
$\xi_i T_i$. Se obtiene, pues, el enunciado siguiente:

\setcounter{subsection}{2}

\begin{corollaire}
\label{XIII.5.3} Sea $X$ un esquema estrictamente local regular de
característica residual $p\geq 0$,
$D=\sum_{1\leq i\leq
\ifthenelse{\boolean{orig}}{n}{r}
} \divisor f_i$ un divisor de cruzamientos normales sobre~$X$,
$U=X-\Supp D$. Pongamos
$$
\tilde U=\varprojlim_{(n_i)}
U[T_1,\dots,T_r]/(T_1^{n_1}-f_1,\dots,T_r^{n_r}-f_r),
$$
donde el límite proyectivo se toma sobre el conjunto ordenado filtrante
(para la relación de divisibilidad) de las familias de enteros $n_i>0$
primos con~$p$. Entonces $\tilde U$ es un recubrimiento universal
moderadamente ramificado de~$U$. Por consiguiente, el grupo fundamental
moderadamente ramificado de~$U$ es
$$
\pi_1^{\tame}(U)\simeq \prod_{\ell\neq p} \ZZ_\ell [1]^r
\qquad \text{(isomorfismo canónico)}
$$
donde se ha puesto
$\ZZ_\ell[1]=\varprojlim_{n>0}\bbmu_{\ell^n}$. El grupo
$\pi_1^{\tame}(U)$ es isomorfo no canónicamente a
$\prod_{\ell\neq p}\ZZ_\ell^r$.
\label{indnot:mm}\oldindexnot{$\ZZ_\ell[1]$|hyperpage}%
\end{corollaire}
